\documentclass{article}

\usepackage{amsmath,amssymb}
\usepackage{xurl}
\usepackage[hidelinks]{hyperref}
\setlength{\emergencystretch}{2em}

% Shared commands for notes in docs/paper-gaps/.

\DeclareUnicodeCharacter{2080}{\ifmmode{}_{0}\else\textsubscript{0}\fi}
\DeclareUnicodeCharacter{2081}{\ifmmode{}_{1}\else\textsubscript{1}\fi}
\DeclareUnicodeCharacter{2082}{\ifmmode{}_{2}\else\textsubscript{2}\fi}
\DeclareUnicodeCharacter{2099}{\ifmmode{}_{n}\else\textsubscript{n}\fi}

\newcommand{\Lean}{\textsc{Lean}}
\ExplSyntaxOn
\NewDocumentCommand{\leanid}{m}{
  \ifmmode
    \textsf{\detokenize{#1}}
  \else
    \group_begin:
    \urlstyle{sf}
    \tl_set:Nn \l_tmpa_tl {#1}
    \tl_replace_all:Nnn \l_tmpa_tl {\_} {_}
    \exp_args:NV \path \l_tmpa_tl
    \group_end:
  \fi
}
\ExplSyntaxOff
\providecommand{\tr}{\operatorname{tr}}
\newcommand{\tnleanIssueBase}{https://github.com/LionSR/TNLean/issues/}
\newcommand{\tnleanPrBase}{https://github.com/LionSR/TNLean/pull/}
\newcommand{\tnleanDocBase}{https://sirui-lu.com/TNLean/docs/}
\newcommand{\ghissue}[1]{\href{\tnleanIssueBase#1}{GitHub issue \##1}}
\newcommand{\ghpr}[1]{\href{\tnleanPrBase#1}{PR \##1}}
\newcommand{\doclink}[1]{\href{\tnleanDocBase#1.html}{\path{#1}}}

% Dirac notation and the identity operator, as in blueprint/src/macros/common.tex.
% \providecommand keeps note-local definitions authoritative.
\usepackage{dsfont}
\providecommand{\ket}[1]{|#1\rangle}
\providecommand{\bra}[1]{\langle #1|}
\providecommand{\braket}[2]{\langle #1 | #2 \rangle}
\providecommand{\ketbra}[2]{|#1\rangle\!\langle #2|}
\providecommand{\Id}{\mathds{1}}
\providecommand{\one}{\mathds{1}}
\providecommand{\C}{\mathbb{C}}
\providecommand{\MN}[1]{M_{#1}(\C)}
\providecommand{\MD}{M_D(\C)}

% External referencing for paper sources.  The build helper writes sanitized
% auxiliary files under build/paper-aux/ and excludes duplicated source labels.
\usepackage{xr}

% Import a generated paper auxiliary file with a caller-supplied label prefix.
% The candidate paths support builds from docs/paper-gaps/, the repository root,
% and build/paper-aux/ itself.
\newcommand{\importpaperaux}[2]{%
  \IfFileExists{../../build/paper-aux/#2.aux}{%
    \externaldocument[#1]{../../build/paper-aux/#2}%
  }{%
    \IfFileExists{build/paper-aux/#2.aux}{%
      \externaldocument[#1]{build/paper-aux/#2}%
    }{%
      \IfFileExists{#2.aux}{%
        \externaldocument[#1]{#2}%
      }{}%
    }%
  }%
}

% General external reference macros
\newcommand{\paperref}[2]{\ref{#1-#2}}
\newcommand{\papereqref}[2]{(\ref{#1-#2})}

% CPSV16 source (1606.00608 / MPDO-22-12-17-2.tex) external document and shortcuts
\importpaperaux{cpsv16-}{1606.00608/MPDO-22-12-17-2-tnlean}
\newcommand{\cpsvref}[1]{\paperref{cpsv16}{#1}}
\newcommand{\cpsveqref}[1]{\papereqref{cpsv16}{#1}}

% Verdict marker: \gapnote{<kind>}{<status>}. Kind is the policy
% classification (clarification, local-correction, scope-restriction,
% unfaithful, false-source, open-gap); status is open, wip (elimination
% underway), resolved, or historical. Machine-read by the site index;
% prints a verdict line.
\newcommand{\gapnote}[2]{\par\noindent\textbf{Verdict:} #1 (#2).\par}


\title{Basis-of-Representatives Scope Restriction in the\\Fundamental Theorem Formalization}
\author{}
\date{2026-05-10}

\begin{document}
\maketitle

\section{The source argument}

Cirac--P\'erez-Garc\'ia--Schuch--Verstraete 2016 states the general
proportional Fundamental Theorem as Corollary~\texttt{II\_cor2} of
\cite{Cirac2016MPDO_arXiv}. The theorem applies to two tensors $A$ and $B$
whose matrix-product vectors are proportional, each presented in its
basis-of-normal-tensor (BNT) canonical form. The canonical form for $A$ is a
direct-sum decomposition
\[
    A^i = \bigoplus_{j=1}^{g_A} I_{r_{A,j}} \otimes \mu_{A,j} (A_{j,1}^i \oplus \cdots \oplus A_{j,r_{A,j}}^i),
\]
where each modulus class $j$ contains $r_{A,j} \ge 1$ representative blocks
$A_{j,1},\ldots,A_{j,r_{A,j}}$ all sharing the same weight modulus
$|\mu_{A,j}|$, and the moduli are strictly ordered across classes:
$|\mu_{A,1}| > |\mu_{A,2}| > \cdots > |\mu_{A,g_A}|$. The same structure
applies to $B$ with multiplicity counts $r_{B,j}$.

The proof of Corollary~\texttt{II\_cor2} in \cite{Cirac2016MPDO_arXiv}
proceeds in several steps. Step~3 (lines~1156--1163 of the arXiv version)
derives the blockwise power-sum identity
\begin{equation}
    \sum_{q=1}^{r_{A,j}} \mu_{A,j,q}^N
    = \sum_{q=1}^{r_{B,j}} (\nu_{B,j,q}\,e^{i\phi_j})^N,
    \qquad \text{for all } j \text{ and all } N \ge 1,
    \label{eq:power-sum-id}
\end{equation}
from the BNT linear-independence input together with the equal-MPV hypothesis
(so the equality of matrix-product vectors forces the equality of the sums of
$N$th powers). Step~4 invokes Lemma~\texttt{Lem:app\_simple} of the same
paper: if two finite sequences $(\lambda_{a,k})_{k=1}^{x_a}$ and
$(\lambda_{b,k})_{k=1}^{x_b}$ satisfy
\[
    \sum_{k=1}^{x_a} \lambda_{a,k}^N = \sum_{k=1}^{x_b} \lambda_{b,k}^N
    \qquad \text{for } N = 1, \ldots, \max(x_a, x_b),
\]
then $x_a = x_b$ and the two sequences are equal as multisets. Applied to
\eqref{eq:power-sum-id}, this recovers $r_{A,j} = r_{B,j} =: r_j$ and the
multiset equality $\{\mu_{A,j,q}\} = \{e^{i\phi_j}\nu_{B,j,q}\}$ for each
modulus class $j$. Step~6 then assembles the global gauge isomorphism as
$Y = \bigoplus_j I_{r_j} \otimes Y_j$, using the per-class dimension matching
$r_{A,j} = r_{B,j}$ from Step~4.

The mathematical structure of this argument is that the multiplicities $r_j$
and the weight multisets are recovered from a finite family of polynomial
identities in $N$, without any passage to a limit.

\section{The basis-of-representatives restriction}

The surviving formal normal-CF-BNT structure\footnote{%
    \path{TNLean/MPS/BNT/Construction.lean}: \leanid{IsNormalCanonicalFormBNT}.
    The earlier non-normal wrapper was retired after its downstream consumers
    were replaced by explicit separated-block hypotheses.} no longer imposes
strictly decreasing moduli. Its weight condition is only non-increasing by
modulus, and all retained weights are nonzero. Thus equal-modulus blocks are
allowed. The additional BNT separation field says instead that distinct blocks
are not gauge-phase equivalent.

The remaining restriction is therefore not strict modulus ordering. It is that
\leanid{IsNormalCanonicalFormBNT} is a basis-of-representatives surface: it
keeps one representative for each gauge-phase class already selected in the
block family. It does not retain, as separate summands of the same structure,
the repeated copies of a single source BNT tensor together with their individual
weights $\mu_{j,q}$, coefficient contribution
\[
    c_N^{(j)}=\sum_q \mu_{j,q}^N,
\]
and copy gauges. That raw multiplicity data is carried in the formalization by
\leanid{SectorDecomposition} and the SectorBNT comparison theorems, not by
\leanid{IsNormalCanonicalFormBNT}.

In the one-copy special case the source decomposition reduces to
$A^i = \bigoplus_{j=1}^{g} \mu_j A_j^i$ with a single copy of each basis tensor.
Then the identity~\eqref{eq:power-sum-id} becomes the single-term equality
$\mu_{A,j}^N = (e^{i\phi_j}\nu_{B,j})^N$, Steps~3--4 of the source proof become
vacuous, and the global gauge in Step~6 takes the simpler form
$Y = \bigoplus_j Y_j$ (no identity tensor factor). The present normal-CF-BNT
surface is a representative-family surface of this kind; the full copy
multiplicity comparison is provided elsewhere.

\section{Missing components for the general case}

Four items are currently absent from the normal-CF-BNT representative surface;
each is required to recover the source-level multiplicity data with
$r_j \ge 1$.

\textbf{Lemma~\texttt{Lem:app\_simple} (finite Newton--Girard uniqueness).}
The source paper states, in its appendix, that two finite sequences whose
power sums agree for $N = 1, \ldots, \max(x_a, x_b)$ must be equal as
multisets. This is a finite algebraic fact, provable by Newton--Girard
identities: the power sums determine the elementary symmetric polynomials,
which determine the multiset of roots. A module named
\leanid{ScalarPowerSumIdentity}\footnote{Introduced in commit
    \texttt{dd54b79e}.} was added at an early stage of the project, but it does
not specialize to this particular fixed-length statement in the form needed by
Step~4 of the source proof. The lemma needs to be formalized as a standalone
result, separate from any BNT context.

\textbf{BNT power-sum identity.}
The identity~\eqref{eq:power-sum-id} requires a proof: starting from the
BNT structure on $A$ and $B$ and the equal-MPV hypothesis, one must derive,
for each modulus class $j$, the equality of the two power sums in
\eqref{eq:power-sum-id}. The formal development no longer packages
coefficient arrays for the discarded proportional branch as standalone
hypotheses. The source derivation of \eqref{eq:power-sum-id} from the BNT
linear-independence data and the equal-MPV condition is still not carried out.
This derivation is logically prior to applying \texttt{Lem:app\_simple}.

\textbf{Multi-copy sector structure.}
The current representative surface \leanid{IsNormalCanonicalFormBNT} allows
equal moduli, but it does not carry repeated gauge-phase-equivalent copies with
their individual weights. The full theorem therefore needs the
\leanid{SectorDecomposition}/SectorBNT multiplicity structure, where each
sector admits an arbitrary finite list of copy weights and gauges. Without
this explicit multiplicity data, the general-multiplicity versions of
Corollary~\texttt{II\_cor2} and of the gauge assembly step cannot be stated on
the normal-CF-BNT surface alone.

\textbf{Global gauge assembly.}
With $r_j \ge 1$ copies per sector, the global gauge isomorphism takes the
form $Y = \bigoplus_j I_{r_j} \otimes Y_j$, where $I_{r_j}$ is the identity
on the multiplicity space and $Y_j$ is the per-block gauge on the bond
dimension. Constructing and verifying this more general gauge requires the
dimension match $r_{A,j} = r_{B,j}$ from Step~4, which in turn requires
\texttt{Lem:app\_simple} and the power-sum identity. The simpler
$Y = \bigoplus_j Y_j$ gauge used in the one-copy case is already present.

\section{What the current formalization does prove}

Within the basis-of-representatives restriction, the formalization establishes
the equal-MPV branch of the Fundamental Theorem argument. The BNT separation
hypotheses yield cross-block overlap decay, and the equal-MPV hypothesis at
a single witness length is used to match the per-block gauge. The former
strict-modulus proportional branch has been deleted: the surviving block-matching
boundary uses the explicit conclusion
\leanid{BlockPermutationGaugePhaseConclusion} only where the BNT comparison
has already been supplied. Thus no current theorem is presented as deriving
the source proportional comparison from the discarded coefficient hypotheses.

The old dominant-normalization field and the old strict-ordering field are no
longer part of \leanid{IsNormalCanonicalFormBNT}. When a CPSV argument needs a
unit-modulus witness for a sector, that datum is carried as an explicit
SectorBNT hypothesis, for example by the corresponding field of
\leanid{SectorDecomposition}.

\section{Downstream audit status}

The downstream CPSV16 surfaces that currently consume the restricted
normal-CF-BNT predicate have been audited against the basis-of-representatives
boundary.\footnote{The relevant Lean surfaces are
    \leanid{IsNormalCanonicalFormBNT},
    \leanid{IsNormalCanonicalFormBNT.ofSeparatedData},
    \leanid{IsNormalCanonicalFormBNT.cross_overlap_tendsto_zero},
    \leanid{IsNormalCanonicalFormBNT.isBNT},
    \leanid{IsNormalCanonicalFormBNT.exists_common_blockTensor_isInjective},
    and
    \leanid{exists_common_blockTensor_isInjective_two_of_isNormalCanonicalFormBNT}.}
Each of these public surfaces either assumes or constructs an already chosen
representative family, non-increasingly ordered by weight modulus. Equal
moduli are allowed, but repeated gauge-phase-equivalent copies and their
individual copy weights are not reconstructed by this surface. These statements
do not introduce a source-faithful canonical-form existence theorem and do not
recover the repeated-copy sector data of the source BNT decomposition. Their
docstrings and the corresponding Wielandt blueprint corollary therefore carry
the same basis-of-representatives marker.

This audit does not apply that marker to the sector-BNT development based on
\leanid{IsBNTCanonicalForm}. That predicate is the paper-faithful sector
surface used for the CPSV16 coefficient and weight-multiplicity arguments,
where repeated copies inside one sector remain explicit.

\section{Conclusion}

The present formalization proves the proportional Fundamental Theorem on the
already selected basis-of-representatives surface. The general case of
Corollary~\texttt{II\_cor2} in \cite{Cirac2016MPDO_arXiv}, which allows
repeated copies inside a sector, requires three additional components not yet
in the formalization: the standalone Newton--Girard uniqueness lemma
(\texttt{Lem:app\_simple}), the derivation of the blockwise power-sum identity
from the BNT structure and the equal-MPV hypothesis, and the global gauge
assembly with multiplicity identity factors. Supplying these components means
threading the existing SectorBNT multiplicity data through the final assembly,
not reviving the deleted strict-modulus route.

The corresponding tracked items are
\href{https://github.com/LionSR/TNLean/issues/1560}{issue~\#1560} for the
standalone \texttt{Lem:app\_simple} lemma,
\href{https://github.com/LionSR/TNLean/issues/1561}{issue~\#1561} for the
BNT power-sum identity derivation, and
\href{https://github.com/LionSR/TNLean/issues/1562}{issue~\#1562} for the
general multi-copy route tracking the full relaxation.

\bibliographystyle{alpha}
\bibliography{references}

\end{document}
